\begin{figure}[t]
\centering
\fbox{
\begin{minipage}{0.95\linewidth}
\textbf{prompt}: You need to come up with the names of some of the papers that relate to the paper you are given.

RoMQA: A Benchmark for Robust, Multi-evidence, Multi-answer Question Answering
\vspace{1em}

\textbf{chosen}:  Sure, here are some papers that relate to RoMDbQA:

1. "Multi-evidence question answering: A survey" by Y. Liu, et al. This paper provides a comprehensive survey of multi-evidence question answering, including its challenges, approaches, and applications. \\
2. "Robust question answering: A study on tolerance to errors in question understanding" by M. G. Rodriguez, et al. This paper investigates the problem of errors in question understanding and proposes a framework for evaluating the robustness of question answering systems. \\
3. "Multi-answer question answering: A review of approaches and challenges" by A. H. Khan, et al. This paper reviews various approaches to multi-answer question answering, discusses the challenges associated with this task, and highlights future research directions. \\
4. "Benchmarking question answering systems: A survey" by A. M. A. Rahim, et al. This paper surveys various benchmarking datasets and evaluation metrics for question answering systems, including those that are specifically designed for multi-evidence and multi-answer question answering. \\
5. "Evaluating the robustness of question answering systems using adversarial attacks" by J. Liu, et al. This paper proposes a method for evaluating the robustness of question answering systems using adversarial attacks and demonstrates its effectiveness on several state-of-the-art systems. \\
6. "Improving question answering robustness using data augmentation" by J. Zhang, et al. This paper investigates the use of data augmentation techniques to improve the robustness of question answering systems and reports promising results on several benchmark datasets. \\
7. "Multi-modal question answering: A review of recent advances" by S. K. Singh, et al. This paper reviews recent advances in multi-modal question answering, which involves answering questions using multiple sources of information, such as text, images, and videos. \\
8. "Question answering in the wild: A study on the robustness of state-of-the-art models" by Y. Zhang, et al. This paper evaluates the robustness of state-of-the-art question answering models on a variety of real-world datasets and discusses the challenges and future research directions in this area.

These papers provide a good starting point for exploring the various aspects of RoMDbQA and related research areas.

\vspace{1em}

\textbf{rejected}:QA is an acronym for Question and Answer. The paper you're asking about is "RoMQA: A Benchmark for Robust, Multi-evidence, Multi-answer Question Answering".

The paper's title is a mouthful, so I'll break it down for you.

1. "RoMQA" stands for "Robust, Multi-evidence, Multi-answer Question Answering".  \\
2. "Benchmark" means that the paper is a standard against which other similar systems can be measured or compared.  \\
3. The paper was written by a team of researchers from the University of Texas at Austin, the University of Washington, and the Chinese University of Hong Kong.  \\

The paper is about a new approach to question answering (QA) that can deal with multiple types of evidence and multiple possible answers. QA is an area of artificial intelligence (AI) that involves developing systems that can answer questions that a human would typically be able to answer.

The RoMQA system is designed to be "robust", meaning that it can still answer questions even when the evidence is noisy or when the question is asked in a way that is different from how it was asked when the system was trained. It can also handle "multi-evidence", meaning that it can use more than one piece of evidence to answer a question, and "multi-answer" questions, which have more than one possible answer.

The RoMQA system was evaluated on a set of questions from a variety of sources, including online forums, news articles, and encyclopedias. The system was able to answer 42\% of the questions correctly, which is an improvement over previous systems.

The paper also describes some of the challenges that still remain in developing effective QA systems, and it outlines directions for future research.

In summary, the paper "RoMQA: A Benchmark for Robust, Multi-evidence, Multi-answer Question Answering" describes a new approach to QA that can handle multiple types of evidence and multiple answers, and it evaluates the system on a set of questions from a variety of sources. The paper also identifies some of the challenges that still remain in QA and outlines directions for future research.

\end{minipage}
}
\caption{Example 1 from the AlpacaEval Length subset, which was rejected due to hallucinations.}
\label{fig:alpacaeval_4}
\end{figure}
